% F -- The two ends of the scaling curve, and the mechanism proposed for each.
\begin{figure}[!htbp]
\centering
\fitwidth{%
\begin{tikzpicture}[font=\small,
  ttl/.style={anchor=north west,align=left,font=\footnotesize\bfseries,text=clmodel},
  axl/.style={font=\tiny,text=clrule},
  gv/.style={font=\scriptsize\bfseries,text=clgain,anchor=north},
]
% ================= panel (a) : the smallest kind becomes estimable ========
\begin{scope}
\node[ttl,text width=8.0cm] at (0,0.40)
  {(a) Why it climbs: the smallest kind of user crosses a hundred members};
\draw[clrule] (0.60,-0.45) -- (0.60,-3.30) -- (8.30,-3.30);
\node[axl,rotate=90,anchor=south] at (0.36,-1.90)
  {people in the evening-moving user};
\draw[clloss,dashed,thick] (0.60,-2.34) -- (8.20,-2.34);
\foreach \xx/\hh/\lbl in {1.55/0.54/56, 3.45/1.21/126, 5.35/1.88/196, 7.25/2.55/266}{
  \fill[clGtwo] (\xx-0.42,-3.30) rectangle (\xx+0.42,{-3.30+\hh});
  \node[font=\scriptsize,text=clrule,anchor=south] at (\xx,{-3.26+\hh}) {\lbl};
}
\node[anchor=north west,align=left,text width=2.10cm,font=\scriptsize,
      text=clloss] at (0.75,-0.62)
  {about 100 members: enough to estimate a rhythm};
\draw[clloss] (1.70,-1.60) -- (1.70,-2.34);
\foreach \xx/\np/\gg in {1.55/400/{+0.6}, 3.45/900/{+2.8}, 5.35/{1{,}400}/{+3.5},
                         7.25/{1{,}900}/{+4.2}}{
  \node[anchor=north,font=\scriptsize,text=clrule] at (\xx,-3.36) {\np};
  \node[gv] at (\xx,-3.78) {\gg};
}
\node[axl,anchor=north] at (4.45,-4.22)
  {top row: people trained on \ \ \ bottom row: what the model gains over the rule};
\end{scope}

% ================= panel (b) : nothing left to add =======================
\begin{scope}[shift={(8.95,0)}]
\node[ttl,text width=5.9cm] at (0,0.40)
  {(b) Why it stops: there is nothing left to see};
\node[anchor=north west,align=left,text width=6.10cm,font=\scriptsize,
      text=clrule] at (0,-0.42)
  {share of the test people's antennas already seen during training};
% bar 1
\node[anchor=west,font=\scriptsize,text=clmodel] at (0,-1.42) {400 people};
\fill[clnonsig!55] (2.60,-1.54) rectangle (5.90,-1.30);
\fill[clbase]      (2.60,-1.54) rectangle (5.86,-1.30);
\node[anchor=west,font=\scriptsize\bfseries,text=clbase] at (6.00,-1.42) {98.9};
% bar 2
\node[anchor=west,font=\scriptsize,text=clmodel] at (0,-1.97) {5{,}000 and above};
\fill[clbase] (2.60,-2.09) rectangle (5.90,-1.85);
\node[anchor=west,font=\scriptsize\bfseries,text=clbase] at (6.00,-1.97) {100};
\draw[clrule!60] (0,-2.38) -- (6.60,-2.38);
\node[anchor=north west,align=left,text width=6.40cm,font=\scriptsize,
      text=clmodel] at (0,-2.52)
  {And the four shapes of Figure~\ref{fig:group-profiles} are essentially
   identical at every size, so 5{,}000 people already describe this population
   completely.};
\node[anchor=north west,draw=clrule,fill=clsoft,rounded corners=2pt,inner sep=5pt,
      align=left,text width=6.05cm,font=\scriptsize] at (0,-3.62)
  {The sixteen thousand people added after that bring no antenna and no
   behaviour the model has not already seen hundreds of times. They are worth
   \textbf{\textcolor{clgain}{+0.1}}.};
\end{scope}
\end{tikzpicture}}
\caption{Why the curve climbs, and why it stops}
\label{fig:why-it-stops}
\figtext{%
\figpara{Panel (a): one mechanism fits the timing of the jump.}{The bars are how
many people fall in the least populated of the four kinds, the evening-moving
user, at 14\,\% of the population. Under each bar is what the model gained over
the one-line rule at that training size. The count crosses roughly a hundred
members exactly between the two sizes where the gain jumps from 0.6 to 2.8. A
label can only carry information if the group it names has enough inhabitants
for its rhythm to be estimated at all, and below that the model is right to
ignore it, which is precisely what the failed run of Chapter~\ref{ch:groups}
showed it doing. Cross-training did not become a better idea at 900 people than
at 400; the kinds simply became large enough to be worth reading.}

\figpara{Panel (b): above the threshold every new person is a duplicate.}{With
113 masts, four kinds of user and every antenna in the test set already seen
during training at 5{,}000 people, more people can only bring more of the same.
This is an explanation that fits, not a proof: the alternative reading, that the
training budget rather than the population ran out, is stated in
Section~\ref{sec:scale-limits}.}}
\figsource{\nbfinal}{train\_cellid\_final.ipynb}{the eight-tier scaling notebook; the member counts are the group shares of Chapter~\ref{ch:groups} applied to each size}
\end{figure}
